\section{A certified exact solver for large instances}
\label{sec:exact}

Excess over the optimum is only meaningful with a trustworthy optimum. At
$N=100{,}000$ a dense matrix does not fit in memory. SciPy's sparse LAPJV on a
16-nearest-cell graph did not finish a 10{,}000-point instance in eight minutes; we
later found that such graphs are usually \emph{infeasible}, because some cells are not
among any point's 16 nearest. We therefore built a solver that never forms the dense
problem, yet certifies optimality against it.

\begin{enumerate}
\item \textbf{Candidates.} Give each point its 12 nearest cells, plus the $5\times5$
  lattice neighbourhood of its cell in a starting assignment. Including the starting
  assignment guarantees that a perfect matching exists.
\item \textbf{Sparse solve.} Solve the min-cost flow (points $\to$ cells $\to$ sink,
  unit capacities) with OR-Tools' cost-scaling solver. The solver needs integer costs,
  so we scale by $S=2^{56}/((N{+}M{+}2)\max c)$ and round.
  (We found by probing that a scale of $2^{58}/((N{+}M)\max c)$ is accepted and
  $2^{60}/((N{+}M)\max c)$ is rejected with \texttt{BAD\_COST\_RANGE}.)
\item \textbf{Potentials.} For the sparse optimum $\sigma$, build a graph on cells with
  an arc $\sigma(i)\to j$ of weight $c_{ij}-c_{i\sigma(i)}$ for every candidate edge.
  Shortest-path distances from a virtual source joined to all cells give potentials
  $v$ that satisfy complementary slackness on every candidate edge, with $v=0$ on empty
  cells. We compute them in exact integer arithmetic with a label-correcting algorithm.
\item \textbf{Global pricing.} Evaluate $\mathrm{LB}(v)$ (Proposition~\ref{prop:lb}) over
  \emph{all} cells. Cells are grouped into $16\times16$ blocks, and a block is scanned
  only if $\mathrm{dist}^2(p_i,\text{block})-\max_{\text{block}}v$ could beat the best
  value found so far. The pass also lists every cell whose reduced cost undercuts the
  point's current cell. If the bound meets the cost, stop. Otherwise add those edges,
  plus the $3\times3$ neighbourhood of each such point's best-priced cell, and go back
  to step 2.
\end{enumerate}

The certificate is the gap $C(\sigma)-\mathrm{LB}(v)$, computed in floating point from
real-valued costs. It is valid whatever errors occur in steps 1--3. Integer rounding
can only loosen it, by at most $4N/S$ in total; we count an instance as certified when
the gap is at most $\max(10^{-9}C, 4N/S)$. The result is therefore the exact optimum of
the dense problem, whatever candidate edges it started from; the start affects only the
time. Wherever SciPy's dense solver fits, both found the same optimum: on 600 small
instances ($N\le2{,}500$, five distributions, three grid types) to relative $10^{-9}$
(Section~\ref{sec:correctness}), and on all ten distributions at $N=10^4$ to relative
$10^{-10}$.

\paragraph{A pitfall.} Our first version flagged a ``negative cycle'' on some clustered
instances with masks. It declared a cycle once a node had been queued more than $n$
times, which is only valid for FIFO Bellman--Ford, not with the small-label-first queue
we use. The fix is to verify a suspected cycle by walking parent pointers. A regression
test covers it.

\paragraph{Cost.} Started from \ShuffleSnap, the solver reached a certificate on every
instance we gave it: the 232 reference instances (up to 200{,}000 points), every timed run,
and uniform instances up to $10^6$ points. Started from bisection, it failed on 4 of 10
instances at $10^5$ points (Section~\ref{sec:main}). The number of pricing rounds was 1 for
uniform inputs, up to about 8 for clustered, curved or masked inputs, and up to 13 for the
input with 16 heavily duplicated sites. Massive ties are the slowest case: one
such 100{,}000-point instance took 34 minutes, while the median 100{,}000-point
reference took under 3 minutes (six references were computed in parallel).
Section~\ref{sec:results} reports its running time and memory next to the heuristics.
The solver needs a feasible starting assignment. We use \ShuffleSnap's default preset,
and at $10^5$ points we also run it from a bisection start, to separate the two methods'
contributions.
